\documentclass{amsart}
% For dealing with references we use the comment environment
\usepackage{verbatim}
\newenvironment{reference}{\comment}{\endcomment}
%\newenvironment{reference}{}{}
\newenvironment{slogan}{\comment}{\endcomment}
\newenvironment{history}{\comment}{\endcomment}

% For commutative diagrams we use Xy-pic
\usepackage[all]{xy}

% We use 2cell for 2-commutative diagrams.
\xyoption{2cell}
\UseAllTwocells

% We use multicol for the list of chapters between chapters
\usepackage{multicol}

% This is generally recommended for better output
\usepackage{lmodern}
\usepackage[T1]{fontenc}

% For cross-file-references

% Package for hypertext links:
\usepackage{hyperref}

% For any local file, say "hello.tex" you want to link to please

% Theorem environments.
%
\theoremstyle{plain}
\newtheorem{theorem}[subsection]{Theorem}
\newtheorem{proposition}[subsection]{Proposition}
\newtheorem{lemma}[subsection]{Lemma}

\theoremstyle{definition}
\newtheorem{definition}[subsection]{Definition}
\newtheorem{example}[subsection]{Example}
\newtheorem{exercise}[subsection]{Exercise}
\newtheorem{situation}[subsection]{Situation}

\theoremstyle{remark}
\newtheorem{remark}[subsection]{Remark}
\newtheorem{remarks}[subsection]{Remarks}

\numberwithin{equation}{subsection}

% Macros
%
\def\lim{\mathop{\mathrm{lim}}\nolimits}
\def\colim{\mathop{\mathrm{colim}}\nolimits}
\def\Spec{\mathop{\mathrm{Spec}}}
\def\Hom{\mathop{\mathrm{Hom}}\nolimits}
\def\Ext{\mathop{\mathrm{Ext}}\nolimits}
\def\SheafHom{\mathop{\mathcal{H}\!\mathit{om}}\nolimits}
\def\SheafExt{\mathop{\mathcal{E}\!\mathit{xt}}\nolimits}
\def\Sch{\mathit{Sch}}
\def\Mor{\mathop{\mathrm{Mor}}\nolimits}
\def\Ob{\mathop{\mathrm{Ob}}\nolimits}
\def\Sh{\mathop{\mathit{Sh}}\nolimits}
\def\NL{\mathop{N\!L}\nolimits}
\def\CH{\mathop{\mathrm{CH}}\nolimits}
\def\proetale{{pro\text{-}\acute{e}tale}}
\def\etale{{\acute{e}tale}}
\def\QCoh{\mathit{QCoh}}
\def\Ker{\mathop{\mathrm{Ker}}}
\def\Im{\mathop{\mathrm{Im}}}
\def\Coker{\mathop{\mathrm{Coker}}}
\def\Coim{\mathop{\mathrm{Coim}}}

% Boxtimes
%
\DeclareMathSymbol{\boxtimes}{\mathbin}{AMSa}{"02}

%
% Macros for moduli stacks/spaces
%
\def\QCohstack{\mathcal{QC}\!\mathit{oh}}
\def\Cohstack{\mathcal{C}\!\mathit{oh}}
\def\Spacesstack{\mathcal{S}\!\mathit{paces}}
\def\Quotfunctor{\mathrm{Quot}}
\def\Hilbfunctor{\mathrm{Hilb}}
\def\Curvesstack{\mathcal{C}\!\mathit{urves}}
\def\Polarizedstack{\mathcal{P}\!\mathit{olarized}}
\def\Complexesstack{\mathcal{C}\!\mathit{omplexes}}
% \Pic is the operator that assigns to X its picard group, usage \Pic(X)
% \Picardstack_{X/B} denotes the Picard stack of X over B
% \Picardfunctor_{X/B} denotes the Picard functor of X over B
\def\Pic{\mathop{\mathrm{Pic}}\nolimits}
\def\Picardstack{\mathcal{P}\!\mathit{ic}}
\def\Picardfunctor{\mathrm{Pic}}
\def\Deformationcategory{\mathcal{D}\!\mathit{ef}}

\setlength{\emergencystretch}{3em}
\begin{document}
\title[Stacks Project, Tag 08XD]{627 --- Descent of module properties along pure ring maps}
\maketitle
\noindent Copyright (C) 2005--2025 Johan de Jong. Stacks Project authors. Source: \href{https://stacks.math.columbia.edu/tag/08XD}{Tag 08XD}.

\noindent Editorial difficulty note (not author text). Difficulty H2: The proof descends finite presentation without flatness and proves flatness using character-module splittings, injectivity, tensor-Hom adjunction, and Tor. This nonflat descent mechanism is substantially beyond ordinary qualification algebra..

\noindent Editorial provenance note (not author text). History: excerpt from revision \texttt{a0444\allowbreak 6e57e\allowbreak c1fbc\allowbreak 252a8\allowbreak 71afc\allowbreak ec775\allowbreak 2fb28\allowbreak 07b14}, descent.tex, lines 1356--1440. Reference presentation was adapted; original-author context and core support blocks were retained or added. The mathematical wording is unchanged. Licensed under GNU FDL 1.2 or later; no invariant sections or cover texts. See ../licenses/COPYING. Context blocks reproduce author definitions and supporting results; they are not additional counted problems.

\begin{definition}
\label{definition-universally-injective}
A ring map $f: R \to S$ is {\it universally injective}
if it is universally injective as a morphism in $\text{Mod}_R$.
\end{definition}

\begin{definition}
\label{definition-C}
Let $R$ be a ring. Define the contravariant functor
{\it $C$} $ : \text{Mod}_R \to \text{Mod}_R$ by setting
$$
C(M) = \Hom_{\textit{Ab}}(M, \mathbf{Q}/\mathbf{Z}),
$$
with the $R$-action on $C(M)$ given by $rf(s) = f(rs)$.
\end{definition}

\begin{remark}
\label{remark-adjunction}
We will use frequently the standard adjunction between $\Hom$ and tensor 
product, in the form of the natural isomorphism of contravariant functors
\begin{equation}
\label{equation-adjunction}
C(\bullet_1 \otimes_R \bullet_2) \cong \Hom_R(\bullet_1, C(\bullet_2)): 
\text{Mod}_R \times \text{Mod}_R \to \text{Mod}_R
\end{equation}
taking $f: M_1 \otimes_R M_2 \to \mathbf{Q}/\mathbf{Z}$ to the map $m_1 \mapsto 
(m_2 \mapsto f(m_1 \otimes m_2))$. See
Algebra, Lemma \href{https://stacks.math.columbia.edu/tag/08YQ}{lemma hom from tensor product variant (Tag 08YQ)}.
A corollary of this observation is that if
$$
\xymatrix@C=9pc{
C(M) \ar@<1ex>[r] \ar@<-1ex>[r] & C(N) \ar[r] & C(P)
}
$$
is a split coequalizer diagram in $\text{Mod}_R$, then so is
$$
\xymatrix@C=9pc{
C(M \otimes_R Q) \ar@<1ex>[r] \ar@<-1ex>[r] & C(N \otimes_R Q) \ar[r] & C(P 
\otimes_R Q)
}
$$
for any $Q \in \text{Mod}_R$.
\end{remark}

Throughout this subsection, fix a universally injective ring map $f : R \to S$,
an object $M \in \text{Mod}_R$, and a ring map $R \to A$. We now investigate
the question of which properties of $M$ or $A$ can be checked after base
extension along $f$. We start with some results from 
\href{https://stacks.math.columbia.edu/bibliography}{Bibliography: mesablishvili2}.


\begin{lemma}
\label{lemma-flat-to-injective}
If $M \in \text{Mod}_R$ is flat, then $C(M)$ is an injective $R$-module.
\end{lemma}

\begin{proof}
Let $0 \to N \to P \to Q \to 0$ be an exact sequence in $\text{Mod}_R$. Since 
$M$ is flat,
$$
0 \to N \otimes_R M \to P \otimes_R M \to Q \otimes_R M \to 0
$$
is exact.
By Lemma \href{https://stacks.math.columbia.edu/tag/08WR}{lemma C is faithful (Tag 08WR)},
$$
0 \to C(Q \otimes_R M) \to C(P \otimes_R M) \to C(N \otimes_R M) \to 0
$$
is exact. By (\ref{equation-adjunction}), this last sequence can be rewritten
as
$$
0 \to \Hom_R(Q, C(M)) \to \Hom_R(P, C(M)) \to \Hom_R(N, C(M)) \to 0.
$$
Hence $C(M)$ is an injective object of $\text{Mod}_R$.
\end{proof}


\begin{lemma}
\label{lemma-split-surjection}
Let $R$ be a ring. A morphism $f: M \to N$ in $\text{Mod}_R$ is universally
injective if and only if $C(f): C(N) \to C(M)$ is a split surjection.
\end{lemma}

\begin{proof}
By (\ref{equation-adjunction}), for any $P \in \text{Mod}_R$ we have a 
commutative diagram
$$
\xymatrix@C=9pc{
\Hom_R( P, C(N)) \ar[r]_{\Hom_R(P,C(f))} \ar[d]^{\cong} &
\Hom_R(P,C(M)) \ar[d]^{\cong} \\
C(P \otimes_R N ) \ar[r]^{C(1_{P} \otimes f)} & C(P \otimes_R M ).
}
$$
If $f$ is universally injective, then $1_{C(M)} \otimes f: C(M) \otimes_R M \to 
C(M) \otimes_R N$ is injective,
so both rows in the above diagram are surjective for $P = C(M)$. We may thus 
lift
$1_{C(M)} \in \Hom_R(C(M), C(M))$ to some $g \in \Hom_R(C(N), C(M))$ splitting 
$C(f)$.
Conversely, if $C(f)$ is a split surjection, then 
both rows in the above diagram are surjective,
so by Lemma \href{https://stacks.math.columbia.edu/tag/08WR}{lemma C is faithful (Tag 08WR)}, $1_{P} \otimes f$ is injective.
\end{proof}


\begin{theorem}
\label{theorem-descend-module-properties}
If $M \otimes_R S$ has one of the following properties as an $S$-module
\begin{enumerate}
\item[(a)]
finitely generated;
\item[(b)]
finitely presented;
\item[(c)]
flat;
\item[(d)]
faithfully flat;
\item[(e)]
finite projective;
\end{enumerate}
then so does $M$ as an $R$-module (and conversely).
\end{theorem}

\begin{proof}
To prove (a), choose a finite set $\{n_i\}$ of generators of $M \otimes_R S$
in $\text{Mod}_S$. Write each $n_i$ as $\sum_j m_{ij} \otimes s_{ij}$ with
$m_{ij} \in M$ and $s_{ij} \in S$. Let $F$ be the finite free $R$-module with
basis $e_{ij}$ and let $F \to M$ be the $R$-module map sending $e_{ij}$ to
$m_{ij}$. Then $F \otimes_R S\to M \otimes_R S$ is surjective, so
$\Coker(F \to M) \otimes_R S$ is zero and hence $\Coker(F \to M)$
is zero. This proves (a).

\medskip\noindent
To see (b) assume $M \otimes_R S$ is finitely presented. Then $M$ is finitely
generated by (a). Choose a surjection $R^{\oplus n} \to M$ with kernel $K$.
Then $K \otimes_R S \to S^{\oplus r} \to M \otimes_R S \to 0$ is exact.
By Algebra, Lemma \href{https://stacks.math.columbia.edu/tag/0519}{lemma extension (Tag 0519)}
the kernel of $S^{\oplus r} \to M \otimes_R S$
is a finite $S$-module. Thus we can find finitely many elements
$k_1, \ldots, k_t \in K$ such that the images of $k_i \otimes 1$ in
$S^{\oplus r}$ generate the kernel of $S^{\oplus r} \to M \otimes_R S$.
Let $K' \subset K$ be the submodule generated by $k_1, \ldots, k_t$.
Then $M' = R^{\oplus r}/K'$ is a finitely presented $R$-module
with a morphism $M' \to M$ such that $M' \otimes_R S \to M \otimes_R S$
is an isomorphism. Thus $M' \cong M$ as desired.

\medskip\noindent
To prove (c), let $0 \to M' \to M'' \to M \to 0$ be a short exact sequence in 
$\text{Mod}_R$. Since $\bullet \otimes_R S$ is a right exact functor,
$M'' \otimes_R S \to M \otimes_R S$ is surjective. So by
Lemma \href{https://stacks.math.columbia.edu/tag/08WR}{lemma C is faithful (Tag 08WR)} the map
$C(M \otimes_R S) \to C(M'' \otimes_R S)$ is injective.
If $M \otimes_R S$ is flat, then
Lemma \ref{lemma-flat-to-injective} shows
$C(M \otimes_R S)$ is an injective object of $\text{Mod}_S$, so the injection
$C(M \otimes_R S) \to C(M'' \otimes_R S)$
is split in $\text{Mod}_S$ and hence also in $\text{Mod}_R$.
Since $C(M \otimes_R S) \to C(M)$ is a split surjection by 
Lemma \ref{lemma-split-surjection}, it follows that 
$C(M) \to C(M'')$ is a split injection in $\text{Mod}_R$. That is, the sequence
$$
0 \to C(M) \to C(M'') \to C(M') \to 0
$$
is split exact. 
For $N \in \text{Mod}_R$, by (\ref{equation-adjunction}) we see that 
$$
0 \to C(M \otimes_R N) \to C(M'' \otimes_R N) \to C(M' \otimes_R N) \to 0
$$
is split exact. By Lemma \href{https://stacks.math.columbia.edu/tag/08WR}{lemma C is faithful (Tag 08WR)}, 
$$
0 \to M' \otimes_R N \to M'' \otimes_R N \to M \otimes_R N \to 0
$$
is exact. This implies $M$ is flat over $R$. Namely, taking
$M'$ a free module surjecting onto $M$ we conclude that
$\text{Tor}_1^R(M, N) = 0$ for all modules $N$ and we can use
Algebra, Lemma \href{https://stacks.math.columbia.edu/tag/00M5}{lemma characterize flat (Tag 00M5)}.
This proves (c).

\medskip\noindent
To deduce (d) from (c), note that if $N \in \text{Mod}_R$ and $M \otimes_R N$ 
is zero,
then $M \otimes_R S \otimes_S (N \otimes_R S) \cong (M \otimes_R N) \otimes_R 
S$ is zero,
so $N \otimes_R S$ is zero and hence $N$ is zero.

\medskip\noindent
To deduce (e) at this point, it suffices to recall that $M$ is finitely 
generated and projective if and only if it is finitely presented and flat.
See Algebra, Lemma \href{https://stacks.math.columbia.edu/tag/00NX}{lemma finite projective (Tag 00NX)}.
\end{proof}
\end{document}
